\chapter{Vom Dichteentwurf zur erneut gerechneten Geometrie}\label{user:topology}
\chaptermark{Topologie: Dichte und Geometrie}
Das vorbereitete Beispiel enthält Gmsh-Netz, zwei normale Model-Dateien und
eine kurze Studien-TOML. Der Quader misst 2 mal 1 mal 1 Meter; die Stirnfläche
ist fest, zwei getrennte Endlasten besitzen je 1 Newton. Das synthetische
Material hat \(E=1000\) Pa und \(\nu=0.3\). Der Beispielaufruf lautet:
\begin{lstlisting}
python -I examples/topology/run.py --output results/topology
\end{lstlisting}
Der neue Ausgabeordner enthält die nachvollziehbaren Eingaben,
\path{result.json}, \path{design.vtu} und Prüfsummen. Im Unterordner
\path{geometry} stehen gültige diskrete Netze sowie akzeptierte und
fehlgeschlagene Reanalysen. Eine bestehende Ausgabe wird nicht überschrieben.
Exitcode 0 verlangt konvergierte Dichteoptimierung und mindestens einen
akzeptierten Geometrieweg. Mehr Iterationen, ein kleineres Residuum oder eine
optisch plausible Form ersetzen keinen erfolgreichen Geometrienachweis.

Die Studien-TOML referenziert beide Model-Dateien mit Gewichten 0.5/0.5,
setzt das Dichtevolumen auf 0.6, den Radius auf 0.3 m und hält passive Zellen
an Last und Lagerung voll. Andere Lastfälle teilen Netz und Material.
Negative Gewichte, unmögliche Volumen, unbekannte Gruppen und bewegte Lager
werden diagnostiziert. Ein größerer Radius glättet den Entwurf stärker und
kann mehr Zwischenmaterial ergeben; er garantiert keine Mindestwandstärke.
Die geometrische 0.9-Volumengrenze ist eine separate Reanalyseentscheidung.
Keine CAD-Rekonstruktion und keine reale Fertigungsfreigabe werden zugesagt.

\section{Rezept topology-run}\label{recipe:topology-run}
\lstinputlisting{examples/topology/run.py}

\section{Rezept topology-prepare}\label{recipe:topology-prepare}
\lstinputlisting{examples/topology/prepare.py}

\section{Rezept topology-verify}\label{recipe:topology-verify}
\lstinputlisting{examples/topology/verify.py}
